\section*{NeurIPS Paper Checklist}

\begin{enumerate}

\item {\bf Claims}
    \item[] Question: Do the main claims made in the abstract and introduction accurately reflect the paper's contributions and scope?
    \item[] Answer: \answerYes{}
    \item[] Justification: The abstract and Section~\ref{sec:intro}
    state the method, theory, scaling, expander, and DiCo claims made
    by the paper. The scope restrictions (IPPO only, VMAS Multi-Agent
    Goal Navigation and Dispersion plus one MPE simple-spread transfer
    panel at $n{=}10$ with categorical policies/TVD; online training at
    $n{=}100$ for PPO and $n{=}50$ for DiCo; $n{=}500$ used for
    frozen-policy timing and expander measurement) are stated in
    Section~\ref{sec:discussion} (Limitations) and do not contradict
    any claim in the abstract or introduction.

\item {\bf Limitations}
    \item[] Question: Does the paper discuss the limitations of the work performed by the authors?
    \item[] Answer: \answerYes{}
    \item[] Justification: Section~\ref{sec:discussion} contains a
    dedicated ``Limitations'' paragraph itemizing five scope limits
    (scale/timing regime coverage, scenario and DiCo-configuration
    coverage, fixed-$G$ ablation, narrow non-Gaussian coverage,
    IPPO-only scope).

\item {\bf Theory assumptions and proofs}
    \item[] Question: For each theoretical result, does the paper provide the full set of assumptions and a complete (and correct) proof?
    \item[] Answer: \answerYes{}
    \item[] Justification: The main-text formal results
    (Propositions~\ref{prop:recovery}, \ref{prop:nonneg},
    \ref{prop:perm}, \ref{prop:complexity},
    \ref{prop:unbiased}, Theorem~\ref{thm:concentration},
    Theorem~\ref{thm:distortion}, and
    Corollary~\ref{cor:expander-distortion}) are proved in full in
    Appendix~\ref{app:proofs}. They are stated with their assumptions in
    Sections~\ref{sec:method}--\ref{sec:theory}; the spectral
    sharpening of Proposition~\ref{prop:spectral-nuclear-discrepancy}
    is stated and proved in Appendix~\ref{app:proofs} with a compact
    summary in Section~\ref{sec:theory}. Assumptions (e.g.,
    non-negative edge weights with $W(G)>0$, bounded distances
    $d(i,j)\in[0,D_{\max}]$, independent Bernoulli inclusion under
    Horvitz--Thompson weights) are stated inside the corresponding
    theorem statements and reused verbatim in the proofs.

\item {\bf Experimental result reproducibility}
    \item[] Question: Does the paper fully disclose all the information needed to reproduce the main experimental results of the paper to the extent that it affects the main claims and/or conclusions of the paper (regardless of whether the code and data are provided or not)?
    \item[] Answer: \answerYes{}
    \item[] Justification: Section~\ref{sec:experiments}
    summarizes the VMAS environment, team sizes, policy architecture,
    training algorithm (IPPO), checkpoints, rollout
    set size, and metric implementation. Appendix~\ref{app:hyperparams}
    lists every PPO hyperparameter and the $n{=}100$ scaling overrides.
    The frozen-init GPU timing sweep and the DiCo integration
    (Section~\ref{sec:experiments:scaling},
    Appendix~\ref{app:verify}, and Section~\ref{sec:experiments:dico-dispersion})
    specify batch
    shape, timing regime (\texttt{torch.cuda.synchronize} on both
    sides of every \texttt{perf\_counter}), number of trials, number
    of seeds, and configuration of the $k$-NN graph and the Bernoulli
    random graph. The $n{=}50$ head-to-head
    (Appendix~\ref{app:dico-n50-sweep}) reuses these hyperparameters
    except for the reduced per-iteration rollout
    (\texttt{on\_policy\_n\_envs\_per\_worker}${=}120$,
    \texttt{on\_policy\_collected\_frames\_per\_batch}${=}12{,}000$)
    and is driven by two matched launchers that share seeds, rollout
    shape, and set points:
    \path{ControllingBehavioralDiversity-fork/scripts/launch_n50_bern_setpoint_sweep.sh}
    for the Bernoulli-$0.1$ arm and
    \path{ControllingBehavioralDiversity-fork/scripts/launch_n50_full_setpoint_sweep.sh}
    for the full-SND arm; the head-to-head comparison, figure, and
    LaTeX table fragment are regenerated by
    \texttt{experiments/n50\_bern\_vs\_full\_comparison.py}. The
    post-hoc complete-graph SND audit in the same appendix is
    regenerated by
    \texttt{experiments/n50\_posthoc\_full\_snd\_validation.py}
    after the corresponding raw GPU logs are available. The
    non-VMAS transfer panel in Appendix~\ref{app:verify} is reproduced by
    \texttt{experiments/mpe\_ippo\_training.py},
    \texttt{experiments/mpe\_measurement\_panel.py}, and
    \texttt{scripts/launch\_mpe\_move2.sh}.
    Everything needed to reconstruct the code paths is
    in \texttt{graphsnd/metrics.py}, \texttt{graphs.py},
    \texttt{training/train\_navigation\_batched.py}, and
    \texttt{experiments/exp2\_timing\_scaling.py} in the companion
    repository.

\item {\bf Open access to data and code}
    \item[] Question: Does the paper provide open access to the data and code, with sufficient instructions to faithfully reproduce the main experimental results, as described in supplemental material?
    \item[] Answer: \answerYes{}
    \item[] Justification: The anonymous supplement contains the code,
    tests, curated result CSV/JSON/PDF artifacts, selected raw DiCo CSV
    logs, and a top-level \texttt{README.md} specifying installation,
    experiment reproduction commands, and a claim-to-code map. A
    preflight script
    (\texttt{scripts/preflight.sh}) runs the test suite and a
    smoke-training pass on each visible CUDA device to catch
    environment problems before overnight runs. For the anonymous
    submission the repository URL is anonymized in the main text; the
    public repository will be released upon acceptance.

\item {\bf Experimental setting/details}
    \item[] Question: Does the paper specify all the training and test details (e.g., data splits, hyperparameters, how they were chosen, type of optimizer) necessary to understand the results?
    \item[] Answer: \answerYes{}
    \item[] Justification: Training details (optimizer, learning
    rate, clip parameter, GAE $\lambda$, discount $\gamma$, rollout
    shape, mini-batch size, PPO epochs, hidden-layer width, activation)
    are listed in Appendix~\ref{app:hyperparams}. The DiCo integration
    of Section~\ref{sec:experiments:dico-dispersion} additionally
    specifies the HetControl architecture, $\SND_{\mathrm{des}}{=}0.1$
    set point, $167$-iteration budget, three seeds, $k{=}3$ for the
    $k$-NN graph, Bernoulli inclusion probability $p{=}0.1$, and the
    $128$-env subsample used for the $k$-NN construction. The
    $n{=}50$ head-to-head
    (Appendix~\ref{app:dico-n50-sweep}) uses the same $167$-iteration
    budget, three seeds per cell, $p{=}0.1$ for the Bernoulli arm,
    and three set points
    $\SND_{\mathrm{des}}\in\{0.12,0.14,0.15\}$; the full-SND arm
    runs the same nine cells with the unmodified DiCo estimator on
    $K_n$ (\texttt{model.diversity\_estimator=full}) for a
    $2{\times}3{\times}3{=}18$-run paired comparison.

\item {\bf Experiment statistical significance}
    \item[] Question: Does the paper report error bars suitably and correctly defined or other appropriate information about the statistical significance of the experiments?
    \item[] Answer: \answerYes{}
    \item[] Justification: The unbiasedness diagnostic
    (Section~\ref{sec:experiments}, Appendix~\ref{app:verify})
    reports $95\%$ confidence intervals from $2{,}000$ Bernoulli draws
    per $p$ and the standardised deviation
    $|\mathrm{bias}|/\mathrm{SE}$ as a sharper test. The concentration
    diagnostic reports the empirical violation rate at $\delta{=}0.1$
    over $2{,}000$ draws per cell and the $95$th-percentile deviation.
    The timing experiments
    (Section~\ref{sec:experiments:scaling},
    Appendix~\ref{app:verify}) report mean, median, and standard
    deviation over $20$ timing trials per cell after $3$ warm-ups.
    The DiCo experiment (Section~\ref{sec:experiments:dico-dispersion})
    reports mean $\pm$ standard deviation over three seeds for reward
    and applied SND, and median with IQR for per-call metric cost.
    The $n{=}50$ head-to-head (Appendix~\ref{app:dico-n50-sweep},
    Table~\ref{tab:dico-n50-bern-vs-full}) reports late-window (last
    $50$ iterations) means $\pm$ cross-seed standard error over three
    seeds per (set-point, estimator) cell and the relative tracking
    error
    $\overline{|\mathrm{applied}\;\SND-\SND_{\mathrm{des}}|}/
    \SND_{\mathrm{des}}$ separately for the Bernoulli-$0.1$ and
    full-SND arms; per-cell and aggregated summary CSVs are
    released alongside the LaTeX fragment. The post-hoc full-SND audit
    reports the same late-window statistics for the complete-graph
    audit signal in
    \texttt{results/dico\_n50\_posthoc\_full\_snd/}.

\item {\bf Experiments compute resources}
    \item[] Question: For each experiment, does the paper provide sufficient information on the computer resources (type of compute workers, memory, time of execution) needed to reproduce the experiments?
    \item[] Answer: \answerYes{}
    \item[] Justification: The $n{\in}\{4,8,16\}$ frozen-checkpoint
    verifications (Section~\ref{sec:experiments},
    Appendix~\ref{app:verify}) run on a single CPU core with the
    reported $20$ trials per cell in minutes. The $n{=}100$ overnight
    PPO run (Section~\ref{sec:experiments:scaling}) fits on one
    RTX~4090, as does the frozen-init $n\in\{100,250,500\}$ GPU timing
    sweep (Section~\ref{sec:experiments:scaling},
    Appendix~\ref{app:verify}). The DiCo
    closed-loop experiment
    (Section~\ref{sec:experiments:dico-dispersion}) also
    runs on one RTX~4090 at the reported $167$-iteration budget.
    The $n{=}50$ head-to-head
    (Appendix~\ref{app:dico-n50-sweep}, $18$ cells: three seeds
    $\times$ three set points $\times$ two estimators) runs on two
    RTX~4090 GPUs in parallel, at ${\approx}2$ GPU-hours per
    Bernoulli-$0.1$ cell ($\approx18$ GPU-hours total) and
    ${\approx}3.5$ GPU-hours per full-SND cell
    ($\approx31.5$ GPU-hours total, dominated by the full-SND
    aggregation over all $\binom{50}{2}{=}1{,}225$ pairs each
    forward). The post-hoc complete-graph audit repeats the same
    $18$-cell grid with an additional logging callback, adding about
    $42$ RTX~4090 GPU-hours. Including this audit, total compute for
    the experiments reported in this paper is under $150$ GPU-hours on
    RTX~4090 class GPUs, plus a few CPU-hours for the
    $n{\in}\{4,8,16\}$ benchmarks and the MPE transfer panel
    (three CPU seeds at $n{=}10$ plus measurement passes).

\item {\bf Code of ethics}
    \item[] Question: Does the research conducted in the paper conform, in every respect, with the NeurIPS Code of Ethics \url{https://neurips.cc/public/EthicsGuidelines}?
    \item[] Answer: \answerYes{}
    \item[] Justification: The work is theoretical and empirical
    research on a metric for behavioral diversity in cooperative
    MARL. It does not involve human subjects, sensitive data, or
    release of models with plausible misuse potential. No deviations
    from the NeurIPS Code of Ethics are required.

\item {\bf Broader impacts}
    \item[] Question: Does the paper discuss both potential positive societal impacts and negative societal impacts of the work performed?
    \item[] Answer: \answerNA{}
    \item[] Justification: Graph-SND is a methodological metric
    contribution that reduces the wall-clock cost of measuring
    behavioral diversity in cooperative MARL. It enables neither
    new capabilities nor new deployments; it only makes an existing
    measurement cheaper. We therefore do not identify a direct path
    to positive or negative societal impact.

\item {\bf Safeguards}
    \item[] Question: Does the paper describe safeguards that have been put in place for responsible release of data or models that have a high risk for misuse (e.g., pre-trained language models, image generators, or scraped datasets)?
    \item[] Answer: \answerNA{}
    \item[] Justification: No data or models with misuse potential
    are released; the released artifacts are experimental code,
    reproducible CSV/JSON result files, and paper figures.

\item {\bf Licenses for existing assets}
    \item[] Question: Are the creators or original owners of assets (e.g., code, data, models), used in the paper, properly credited and are the license and terms of use explicitly mentioned and properly respected?
    \item[] Answer: \answerYes{}
    \item[] Justification: VMAS \citep{bettini2022vmas} and DiCo
    \citep{bettini2024dico} are credited in the main text and
    released under their original open-source licenses. The vendored
    fork of the DiCo codebase in
    \texttt{ControllingBehavioralDiversity-fork/} retains the
    upstream license, adds a \texttt{GRAPH\_SND\_CHANGES.md} file
    documenting the modifications, and includes the preferred
    citation in \texttt{CITATION.cff}. The original SND paper
    \citep{bettini2025snd} is cited throughout.

\item {\bf New assets}
    \item[] Question: Are new assets introduced in the paper well documented and is the documentation provided alongside the assets?
    \item[] Answer: \answerYes{}
    \item[] Justification: The new assets are the \texttt{graphsnd}
    Python package, the two training scripts, the experiment
    scripts, and the released figure and CSV outputs. The top-level
    \texttt{README.md} documents installation, the four entry-point
    experiments, the overnight training launcher, and the repository
    layout, and cross-references the specific paper sections that
    each asset reproduces.

\item {\bf Crowdsourcing and research with human subjects}
    \item[] Question: For crowdsourcing experiments and research with human subjects, does the paper include the full text of instructions given to participants and screenshots, if applicable, as well as details about compensation (if any)? 
    \item[] Answer: \answerNA{}
    \item[] Justification: No crowdsourcing or human-subjects
    research is involved.

\item {\bf Institutional review board (IRB) approvals or equivalent for research with human subjects}
    \item[] Question: Does the paper describe potential risks incurred by study participants, whether such risks were disclosed to the subjects, and whether Institutional Review Board (IRB) approvals (or an equivalent approval/review based on the requirements of your country or institution) were obtained?
    \item[] Answer: \answerNA{}
    \item[] Justification: No human-subjects research is involved.

\item {\bf Declaration of LLM usage}
    \item[] Question: Does the paper describe the usage of LLMs if it is an important, original, or non-standard component of the core methods in this research? Note that if the LLM is used only for writing, editing, or formatting purposes and does \emph{not} impact the core methodology, scientific rigor, or originality of the research, declaration is not required.
    \item[] Answer: \answerNA{}
    \item[] Justification: No LLM is part of the methods, theory,
    or experimental pipeline of this paper.

\end{enumerate}
